% ID: SEL-D0031-OB-E02
% ESTADO: validada. S032, cotejada con examen y criterios.
\textbf{a)} Los tramos son continuos en sus respectivos
dominios, pues $25+3x>0$ para $x>50$. En $50$, el primer
tramo vale $400/100=4$ y el límite derecho es $700/175=4$.
Hay también continuidad por la derecha en $10$. En conjunto
$C$ es continua en $[10,\infty)$.

\textbf{b)} En el primer tramo la pendiente es $5/100>0$.
En el segundo, la derivada por la regla del cociente es
\[
C'(x)=\frac{10(25+3x)-3(200+10x)}{(25+3x)^2}
=\frac{-350}{(25+3x)^2}<0.
\]
La cantidad dedicada a créditos decrece a partir de una liquidez
de 50000 euros. El máximo es $C(50)=4$ miles de euros.

\textbf{c)} Al dividir numerador y denominador por $x$ resulta
$\lim_{x\to\infty}C(x)=10/3$. La asíntota horizontal es $C=10/3$.
Para una liquidez cada vez mayor, la cantidad destinada a créditos
se aproxima desde arriba a unos 3333.33 euros.
